%%%%%%%%%%%%%%%%%%%%%%%%%%%%%%%%%%%%%%%%%%%%%%%%%%%%%%%%%%%%%%%%%%%%%%
% Lớp 11 | Đường thẳng, mặt phẳng. Quan hệ song song trong | Điểm, đường thẳng và mặt phẳng
% Dạng 1H4?1-7: Toán thực tế áp dụng điểm, đường thẳng và mặt phẳng
%%%%%%%%%%%%%%%%%%%%%%%%%%%%%%%%%%%%%%%%%%%%%%%%%%%%%%%%%%%%%%%%%%%%%%


% nguồn: 2026-11-15_11_1H4_TL2_Song_song (nhập 15/11/2026)
\begin{bt}%[1H4C1-7]%[ma:58d69635]
% nguon: Archive/SP_KhoLatex/ID6_old/1-CK1-OT-De-19.tex #2877
%[Dự án EX-CK1 (10,11 mới)- Lâm Chính]
	Một khối gỗ hình hộp có tất cả các mặt là các hình vuông, người thợ mộc muốn cắt khối gỗ theo một mặt phẳng sao cho vết cắt trên các mặt của khối gỗ tạo thành một đa giác đều không phải hình vuông (không nhất thiết phải cắt tất cả các mặt của khối gỗ). Tính diện tích lớn nhất của đa giác đều đó khi cạnh của hình vuông bằng $1$ (m).
	\loigiai{
		\begin{itemize}
			\item \textbf{Trường hợp 1.} Đa giác là tam giác đều.\\
			Khi đó cạnh của đa giác đều nằm trọn vẹn trong hình vuông nên diện tích tam giác đều lớn nhất khi cạnh của nó lớn nhất bằng $\sqrt{2}$ (m) từ đó diện tích tam giác đều lớn nhất bằng $\dfrac{\sqrt{3}}{2} \left(m^2 \right)$.
			\item \textbf{Trường hợp 2.} Đa giác là ngũ giác đều.\\
			Khi đó sẽ có 5 mặt của khối gỗ được cắt, nên đa giác đều sẽ phải có 2 cặp cạnh song song. Trường hợp này không xảy ra vì ngũ giác đều không có cặp cạnh nào song song.
			\immini{
			\item \textbf{Trường hợp 3.} Đa giác là lục giác đều.\\
			Ta có $MNPQEF$ là đa giác đều tâm $O$.\\
			 Ba mặt phẳng $\left(MNPQEF\right),\left(BDD'B'\right),\left(ABCD\right)$ cắt nhau theo 3 giao tuyến $MN,BD,FP$ đôi một song song (vì $MN  \parallel  FP$).\\
			Suy ra $MN,BD,FP,EQ,B'D'$ song song với nhau.\\
			Mặt khác $\dfrac{DM}{AM}=\dfrac{BN}{AN}=\dfrac{BP}{PB'}=\dfrac{QC'}{QB'}$.\\
			$\Rightarrow MQ  \parallel  AB'\Rightarrow MQ=AB'=\sqrt{2} \Rightarrow MN=\dfrac{\sqrt{2}}{2}$.\\
			Suy ra $S_{MNPQEF}=6\cdot S_{OMN}=\dfrac{3\sqrt{3}}{4} \left(m^2 \right)$.\\
			Vậy diện tích lớn nhất bằng $\dfrac{3\sqrt{3}}{4} \left(m^2 \right)$.}
			{
			\begin{tikzpicture}[scale=0.8, font=\footnotesize,line join=round, line cap=round, >=stealth]
				\coordinate (A) at (0,0);
				\coordinate (D) at (-2,-2);
				\coordinate (B) at (5,0);
				\coordinate (C) at ($(B)+(D)-(A)$);
				\foreach \i in {A,B,C,D}{\coordinate (\i') at ($(\i)-(0,4)$);}
				\coordinate (M) at ($(A)!0.5!(D)$);
				\coordinate (N) at ($(A)!0.5!(B)$);
				\coordinate (P) at ($(B)!0.5!(B')$);
				\coordinate (Q) at ($(B')!0.5!(C')$);
				\coordinate (E) at ($(C')!0.5!(D')$);
				\coordinate (F) at ($(D')!0.5!(D)$);
				\coordinate (O) at ($(M)!0.5!(Q)$);
				\draw (A)--(B)--(C)--(D)--cycle;
				\draw (B)--(B') (C)--(C') (D)--(D')  (B')--(C')--(D') (M)--(N) (P)--(Q) (E)--(F);
				\draw[dashed,thin](B')--(A')--(A) (A')--(D') (E)--(N)--(P)--(F)--(M)--(Q)--(E);
				\foreach \i/\g in {A'/-90,B'/-90,C'/-90,D'/-90,A/90,B/90,C/90,D/90,M/90,N/90,P/0,Q/-90,E/-90,F/180,O/45}{\draw[fill=black](\i) circle (1.5pt) ($(\i)+(\g:3mm)$) node[scale=1]{$\i$};}
			\end{tikzpicture}
			}
		\end{itemize}

	}
\end{bt}

% nguồn: 2026-11-15_11_1H4_TL2_Song_song (nhập 15/11/2026)
\begin{bt}%[1H4H1-7]%[ma:480114be]
% nguon: Archive/SP_KhoLatex/ID6_old/1-GK1-CanhDieu-De6-NH23-24.tex #3025

	Trên một mảnh đất hình vuông $ A B C D $, bác An đặt một chiếc đèn pin tại vị trí $ A $ chiếu chùm sáng phân kì sang phía góc $ C $. Bác An nhận thấy góc chiếu sáng của đèn pin giới hạn bởi hai tia $ A M, A N $, ở đó các điểm $ M, N $ lần lượt thuộc các cạnh $ B C, C D $ sao cho $ B M=\dfrac{1}{2} B C, D N=\dfrac{1}{2} D C$.
	\begin{enumerate}
		\item Tính $ \tan (B A M+D A N) $. 
		\item Góc chiếu sáng của đèn pin bao nhiêu độ? 
	\end{enumerate}
	\loigiai{
		\immini{
				\begin{enumerate}
				\item Trong tam giác vuông $ A B M $, có $ \tan B A M=\dfrac{B M}{B A}=\dfrac{1}{2} $.\\ Trong tam giác vuông $ A D N $, có $ \tan D A N=\dfrac{D N}{A D}=\dfrac{D N}{D C}=\dfrac{1}{3} $.\\
				Do đó, $$ \tan (B A M+D A N)=\dfrac{\tan B A M+\tan D A N}{1-\tan B A M \cdot \tan D A N}=\dfrac{\dfrac{1}{2}+\dfrac{1}{3}}{1-\dfrac{1}{2} \cdot \dfrac{1}{3}}=1.$$ 
				\item Từ câu a) ta có $ \tan (B A M+D A N)=1 $ nên $ B A M+D A N=45^{\circ} $.\\ Suy ra $ M A N=B A D-(B A M+D A N)=90^{\circ}-45^{\circ}=45^{\circ} $. \\
				Vậy góc chiếu sáng của đèn pin là $ 45^{\circ} $. 
			\end{enumerate}
		}{
			\begin{tikzpicture}
				\coordinate (A) at (0,0);
				\coordinate (B) at (3,0);
				\coordinate (D) at (0,3);
				\coordinate (C) at (3,3);
				\coordinate (N) at ($(D)!0.5!(C)$);
				\coordinate (M) at ($(B)!0.5!(C)$);
				\draw (A)--(B)--(C)--(D)--(A) (A)--(M) (A)--(N);
				\foreach \x/\y in {A/-90,B/-90,C/90,D/90, M/0, N/90}{\fill (\x) circle(1pt) ($(\x)+(\y:0.3cm)$)node{$\x$};}
			\end{tikzpicture}
		}
	
		}
	
\end{bt}
